% ============================================================
% แผนการสอนและประมวลรายวิชา (Course Syllabus)
% ============================================================
\chapter*{ประมวลรายวิชาและแผนการจัดการเรียนรู้}
\addcontentsline{toc}{chapter}{ประมวลรายวิชาและแผนการจัดการเรียนรู้}

\noindent\textbf{รหัสและชื่อวิชา} \quad AIOT4302 ระบบสื่อสารไร้สายและอินเทอร์เน็ตของสรรพสิ่งทางการเกษตร\\
(Wireless Communications and Agricultural Internet of Things)

\noindent\textbf{จำนวนหน่วยกิต} \quad 3(2-2-5) บรรยาย 2 ชั่วโมง ปฏิบัติการ 2 ชั่วโมง ศึกษาค้นคว้าด้วยตนเอง 5 ชั่วโมงต่อสัปดาห์

\noindent\textbf{คำอธิบายรายวิชา} \quad หลักการและสถาปัตยกรรมระบบสื่อสารไร้สายระยะสั้นและระยะปานกลางสำหรับงานเกษตรกรรมแม่นยำสูง โพรโทคอลเครือข่ายแบบเมชตามมาตรฐาน IEEE 802.15.4 และ Zigbee 3.0 สถาปัตยกรรมการจัดเส้นทางแบบพลวัต ระบบนิเวศเซนเซอร์ตรวจวัดดิน บรรยากาศ และรังสีแสง การเชื่อมโยงเซนเซอร์อุตสาหกรรมผ่านสะพานสื่อสาร RS485 Modbus RTU สถาปัตยกรรมเกตเวย์โครงข่ายร่วมกับ Zigbee2MQTT และโพรโทคอล MQTT การพัฒนาเฟิร์มแวร์ไมโครคอนโทรลเลอร์ ESP32-C6 สถาปัตยกรรม RISC-V การคำนวณและบริหารจัดการพลังงานต่ำ และการออกแบบโครงงานบูรณาการในแปลงเกษตรอัจฉริยะ

\vspace{0.4cm}
\noindent\textbf{แผนการจัดการเรียนรู้รายสัปดาห์}

\begin{table}[H]
\caption{แผนการจัดการเรียนรู้รายสัปดาห์ตลอดภาคการศึกษา}
\label{tab:syllabus_schedule}
\centering
\small
\begin{tabularx}{\textwidth}{c p{3.2cm} >{\raggedright\arraybackslash}X c}
\toprule
\textbf{สัปดาห์} & \textbf{หัวข้อเนื้อหา} & \textbf{กิจกรรมและการทดลองปฏิบัติการ} & \textbf{สื่อและอุปกรณ์} \\
\midrule
1 & รากฐานฟิสิกส์คลื่นวิทยุ & วิเคราะห์การลดทอนสัญญาณคลื่น 2.4 GHz ในแปลงเกษตร & เครื่องวิเคราะห์สเปกตรัม \\
2--3 & สถาปัตยกรรม Zigbee 3.0 & การกำหนดบทบาท Coordinator, Router และ End Device & ซอฟต์แวร์ Wireshark \\
4 & อัลกอริทึม AODV Routing & จำลองการจัดเส้นทางแบบเมชและการกู้คืนเครือข่าย & บอร์ดทดลอง CC2652P \\
5 & เปรียบเทียบมาตรฐานไร้สาย & ทดสอบสมรรถนะ Zigbee เทียบกับ LoRaWAN และ Wi-Fi & ชุดรับส่งสัญญาณทดสอบ \\
6--7 & เซนเซอร์ดินอัจฉริยะ & ติดตั้งโพรบวัดความชื้น อุณหภูมิ EC, pH และธาตุอาหาร NPK & โพรบดิน RS485 Modbus \\
8 & เซนเซอร์บรรยากาศและรังสี & บูรณาการหัววัด SHT30, BMP280 และโฟโตเซนเซอร์ PAR & สถานีสภาพอากาศจำลอง \\
9 & สะพานเชื่อมต่อ RS485 & พัฒนาเฟิร์มแวร์แปลงค่า Modbus RTU สู่โพรโทคอล Zigbee & โมดูล MAX485 และ MCU \\
10 & สอบวัดผลกลางภาค & ประเมินผลภาคทฤษฎีและทดสอบทักษะโครงข่ายเชิงปฏิบัติการ & ชุดข้อสอบและใบงาน \\
11 & เกตเวย์ Zigbee2MQTT & ติดตั้งและตั้งค่า Coordinator CC2652P บน Raspberry Pi & เกตเวย์คอมพิวเตอร์ \\
12 & โครงสร้างข้อมูล MQTT & กำหนดสเปกหัวข้อ JSON Telemetry และคำสั่งควบคุม Actuator & ซอฟต์แวร์ MQTT Explorer \\
13 & สถาปัตยกรรม ESP32-C6 & พัฒนาเฟิร์มแวร์ด้วย ESP-IDF และ ESPHome บนคลื่น 802.15.4 & บอร์ด ESP32-C6 DevKit \\
14 & การจัดการพลังงานต่ำ & ทดลอง Deep Sleep คำนวณแบตเตอรี่ และออกแบบโซลาร์เซลล์ & เครื่องวัดกระแสละเอียด \\
15 & โครงงานบูรณาการ & นำเสนอผลงานแปลงทดลองเสมือนจริง 120 โหนด และประเมินผล & แปลงทดสอบภาคสนาม \\
\bottomrule
\end{tabularx}
\end{table}

\clearpage
